\documentclass[12pt,a4paper]{report}

% --- PACKAGES ESSENTIELS ---
\usepackage[french]{babel}
\usepackage[T1]{fontenc}
\usepackage[utf8]{inputenc}
\usepackage{amsmath}
\usepackage{amsfonts}
\usepackage{listings}
\usepackage{xcolor}
\usepackage[most]{tcolorbox}
\usepackage{geometry}

\geometry{margin=2.5cm}

% --- CONFIGURATION DES COULEURS (Style VS Code) ---
\definecolor{codegreen}{rgb}{0,0.6,0}
\definecolor{codegray}{rgb}{0.5,0.5,0.5}
\definecolor{codepurple}{rgb}{0.58,0,0.82}
\definecolor{backcolour}{rgb}{0.95,0.95,0.92}
\definecolor{titleblue}{rgb}{0.1, 0.3, 0.6}

% --- STYLE DES BOITES DE CODE (Correction du chevauchement) ---
\newtcolorbox{codeblock}[1]{
    enhanced,
    attach boxed title to top left={yshift=-2mm, xshift=2mm},
    colback=white,
    colframe=titleblue,
    fonttitle=\bfseries,
    colbacktitle=titleblue,
    title=#1,
    boxed title style={size=small, sharp corners},
    drop shadow
}

% --- CONFIGURATION LISTINGS ---
\lstset{
    language=C,
    backgroundcolor=\color{backcolour},   
    commentstyle=\color{codegreen},
    keywordstyle=\color{magenta},
    numberstyle=\tiny\color{codegray},
    stringstyle=\color{codepurple},
    basicstyle=\ttfamily\footnotesize,
    breaklines=true,                 
    showstringspaces=false,
    tabsize=4
}

\begin{document}

% --- PAGE DE GARDE ---
\begin{titlepage}
    \centering
    \vspace*{2cm}
    {\Huge\bfseries Documentation Technique : Widgets GTK4\par}
    \vspace{1cm}
    {\Large Module de Gestion des Boutons\par}
    \vfill
    {\Large Présenté par : \textbf{Yahya}\par}
    {\large Étudiant en première année de génie logiciel\par}
    \vfill
    {\large \today \par}
\end{titlepage}

\tableofcontents
\newpage

\chapter{Définitions et Structures (Header)}

\section{Configuration des Boutons Standards}

\begin{codeblock}{Structure : button\_config}
La structure \texttt{button\_config} permet de centraliser tous les paramètres d'un bouton (label, icône, style, et signaux) en un seul objet. Cela rend le code plus modulaire.
\end{codeblock}

\begin{lstlisting}
typedef struct {
    const char *label;
    const char *icon_name;
    const char *theme_variant;
    bool use_underline;
    bool has_frame;
    bool can_shrink;
    int icon_size;
    widget_style_config style;
    void (*on_clicked)(GtkButton *button, gpointer user_data);
    gpointer user_data;
} button_config;
\end{lstlisting}

\section{Macros d'Initialisation par Défaut}

\begin{codeblock}{Macro : BUTTON\_CONFIG\_DEFAULT}
Cette macro est cruciale pour éviter les comportements indéfinis. Elle garantit que chaque champ est initialisé à une valeur sûre (souvent NULL ou false) avant d'être modifié.
\end{codeblock}

\begin{lstlisting}
#define BUTTON_CONFIG_DEFAULT ((button_config){ \
    .label = NULL, \
    .icon_name = NULL, \
    .theme_variant = NULL, \
    .use_underline = false, \
    .has_frame = true, \
    .can_shrink = true, \
    .icon_size = 0, \
    .style = WIDGET_STYLE_CONFIG_DEFAULT, \
    .on_clicked = NULL, \
    .user_data = NULL, \
})
\end{lstlisting}

\chapter{Implémentation (Fichier Source)}

\section{Gestion des Thèmes}

L'application utilise des classes CSS personnalisées pour gérer l'apparence. La fonction statique suivante permet d'appliquer ces variantes.

\begin{codeblock}{Fonction : apply\_theme\_variant}
Elle vérifie la chaîne de caractères passée en paramètre et ajoute la classe CSS correspondante au widget GTK.
\end{codeblock}

\begin{lstlisting}
static void apply_theme_variant(GtkWidget *button, const char *variant) {
    if (variant == NULL || variant[0] == '\0') return;

    if (strcmp(variant, "primary") == 0) {
        gtk_widget_add_css_class(button, "btn-primary");
    } else if (strcmp(variant, "success") == 0) {
        gtk_widget_add_css_class(button, "btn-success");
    } else if (strcmp(variant, "danger") == 0) {
        gtk_widget_add_css_class(button, "btn-danger");
    }
}
\end{lstlisting}

\section{Création de Widgets}

\begin{codeblock}{Fonction : create\_button}
C'est le cœur du module. Elle instancie un \texttt{GtkButton}, configure son apparence, ajoute éventuellement une icône via une \texttt{GtkBox}, et connecte le signal "clicked" à la fonction de rappel (callback) fournie.
\end{codeblock}

\begin{lstlisting}
GtkWidget *create_button(const button_config *config) {
    GtkWidget *button;
    // ... Logique d'initialisation GTK4 ...
    if (config->on_clicked) {
        g_signal_connect(button, "clicked", G_CALLBACK(config->on_clicked), config->user_data);
    }
    return button;
}
\end{lstlisting}

\end{document}